
\section{Extended methods}


\subsection{Pre-processing of protein structure data}
To pre-process protein structure data, we developed two distinct pipelines based on either (i)
PyRosetta~\citep{chaudhury_pyrosetta_2010} or (ii) Biopython~\citep{cock_biopython_2009} together
with other open-source tools, with code adapted from~\citep{blaabjerg_rapid_2023}. The
PyRosetta-based pipeline is considerably faster but requires a license, whereas the Biopython-based
pipeline is fully open source. We train models using data generated by both pipelines, with the
constraint that the Pre-processing pipeline used at inference must match that used during training.
Performance differences between the two pipelines are minor (Fig.~\ref{fig:rasp_exp}
and~\ref{fig:t2837_py_vs_bp}, and Table~\ref{table:aa_wt_cls}); unless otherwise stated, we report
results obtained with the PyRosetta pipeline, which yields slightly better performance and has
faster Pre-processing runtime.

For the PyRosetta workflow, we use PyRosetta functionalities for all Pre-processing steps: repairing
PDB files by adding missing residues, adding hydrogen atoms, assigning partial atomic charges,
computing solvent-accessible surface areas (SASA); we ignore non-canonical amino acids.
For the open-source Pre-processing workflow (``Biopython" pipeline), we proceed as follows. First,
we use OpenMM~\citep{eastman_openmm_2013} to repair PDB files by adding missing residues and
substituting non-canonical residues with their canonical counterparts. Hydrogen atoms are then added
using the \texttt{reduce} program~\citep{word_asparagine_1999}. Partial atomic charges are assigned
from the AMBER99sb force field~\citep{ponder_force_2003}, and SASA are computed using
Biopython~\citep{cock_biopython_2009}.

Both pre-processing pipelines retain atoms belonging to non-protein residues and ions, in contrast
to the RaSP pre-processing procedure~\citep{blaabjerg_rapid_2023}. Notably, the PyRosetta-based
pipeline does not replace non-canonical residues.


\subsection{HERMES pre-training}
We pre-trained HERMES using an inverse folding objective, in which the model predicts the identity
of a masked focal amino acid (i.e., the native residue in the structure) given its surrounding
atomic neighborhood. This training task is analogous to that used for
H-CNN~\citep{pun_learning_2024}. We adopted the same data splits as in H-CNN: 10,957 structures for
training, 2,730 for validation, and 212 for testing. These splits are derived from
ProteinNet's~\citep{alquraishi_proteinnet_2019} 30\% sequence-identity clustering of PDB structures
available at the time of CASP12.

Model parameters were optimized for 10 epochs using the Adam optimizer~\citep{kingma_adam_2015} with
a learning rate of $10^{-3}$. We selected the model checkpoint with the lowest validation loss at
the end of each epoch. A single HERMES model is implemented as an ensemble of 10 independently
trained neural network instances, whose predictions are averaged at inference time. Pre-training a
single network instance required approximately 40 minutes per epoch on a single NVIDIA A40 GPU.


\subsection{Rosetta relaxation of mutant structures for the HERMES-{\em relaxed} protocol}
For the HERMES-{\em relaxed} protocol, we use PyRosetta~\citep{chaudhury_pyrosetta_2010} to generate
and locally refine mutant protein structures. Specifically, the focal residue is first substituted
with the desired mutant amino acid, after which all side-chain atoms within 12~\AA~of the focal
residue's C-$\alpha$ atom are relaxed using the FastRelax protocol with the \texttt{ref2015\_cart}
energy function and a single relaxation cycle.

We performed a targeted ablation study to select the relaxation parameters. These experiments
indicate that allowing backbone atoms to relax slightly degrades performance, and that ensembling
over multiple independent relaxation runs does not improve accuracy
(Fig.~\ref{fig:relaxations_ablation}A), while substantially increasing computational cost
(Fig.~\ref{fig:relaxations_ablation}B). Thus, all results for the HERMES-{\em relaxed} protocol are
reported using side-chain–only relaxation with a single FastRelax cycle.

\subsection{Training HERMES-\emph{amortized} to implicitly learn about relaxation}
We fine-tuned HERMES to align HERMES-\emph{fixed} predictions to those of HERMES-\emph{relaxed}, on
a subset of the pre-training protein sites. Specifically, we uniformly sampled 10\% of the proteins
pre-training training set (1,284), and from those we uniformly sampled 5\% of the sites. We
similarly sampled 1\% of the sites from the proteins of the pre-training validation set, and 5\% of
the pre-training testing set.

\subsection{HERMES fine-tuning for downstream tasks}
\label{sec:finetuning_details}
In all analyses, we fine-tuned HERMES models under a Huber Loss objective with hyperparameter
$\delta = 1.0$, for 15 epochs using the Adam optimizer~\citep{kingma_adam_2015} with a learning rate
of $10^{-3}$ and a batch size of 128 mutations, selecting the model checkpoint with the lowest
validation loss at the end of each epoch. The only exception is the models without pre-training,
which are trained for 25 epochs.

By convention, we define model outputs such that higher predicted values correspond to more
favorable mutation effects. Accordingly, when fine-tuning on experimental $\Delta\Delta G$
measurements--where lower values indicate greater stability--we trained the model to predict
$-\Delta\Delta G$. This sign convention ensures consistent interpretation of model outputs across
tasks and is fixed throughout all reported analyses.

To speed-up convergence during fine-tuning, we first rescale the weight matrix and bias vector of
the network's output layer so that the mean and variance of the output logits match those of the
validation scores. This initialization step requires a forward pass through the validation data to
estimate the mean and variance, but it makes the model outputs immediately on the same scale as the
experimental scores. Thus, fine-tuning avoids spending early epochs merely adjusting output
magnitudes.

Overall, fine-tuning is computationally efficient. Training a single neural network instance
requires approximately 2.5 minutes per epoch on the cDNA117k dataset ($\sim$117,000 mutations) and
about 4 minutes per epoch on the Megascale training dataset ($\sim$217,000 mutations) on a single
NVIDIA A40 GPU.






\subsection{Fine-tuning datasets}
\label{sec:datasets}


\noindent {\bf Stability effect prediction.} For protein stability prediction, we fine-tuned and
evaluated the performance of HERMES on the same datasets of three recent structure-based predictors
of mutational stability effects--RaSP~\citep{blaabjerg_rapid_2023},
Stability-Oracle~\citep{diaz_stability_2024}, and ThermoMPNN~\citep{dieckhaus_transfer_2024}--so
results are comparable. The data from these models are used in following way:



\begin{itemize}
\item {\em RaSP dataset.} We used the RaSP dataset~\citep{blaabjerg_rapid_2023} as provided by the
authors on their github repository (\url{https://github.com/KULL-Centre/_2022_ML-ddG-Blaabjerg}). As
in the original RaSP paper, the target values are Rosetta-computed stability changes ($\Delta\Delta
G$), which are known to be reliable primarily within the range [-7, 1] kcal/mol. To account for
this, RaSP applies a sigmoid transformation to the raw $\Delta\Delta G$ values prior to training,
effectively saturating the targets outside this interval.

We adopt the same sigmoid transformation, with one modification: we center the transformed values
such that $\Delta\Delta G = 0$ maps to zero after transformation. This centering is required by the
HERMES output parameterization, in which the predicted stability change for a mutation to the same
amino acid is constrained to be zero (i.e., $\Delta\Delta G_{aa_i \rightarrow aa_i} = 0$). While
this property holds for physical $\Delta\Delta G$ values, it is not preserved by the uncentered
sigmoid transform used in RaSP. Our centered transformation therefore ensures consistency between
the model's output space and the physical interpretation of stability changes, and it follows,

\begin{equation}
    F(\Delta\Delta G) = \frac{1}{1 + e^{-\beta(\Delta\Delta G - \alpha)}} - \frac{1}{1 + e^{\beta
    \alpha}}
    \label{eq:centered_fermi_transform}
\end{equation}\\

\item {\em Stability-Oracle datasets.} Stability-Oracle introduces two curated datasets that we use
in this work~\citep{diaz_stability_2024}:\\

(i) {cDNA117k} (training set): derived from the Megascale cDNA display proteolysis dataset \#
1~\citep{tsuboyama_mega-scale_2023}. The original Megascale dataset reports approximately 850,000
thermodynamic folding stability measurements ($\Delta G$) across 354 natural and 188 de novo
mini-protein domains (40–72 amino acids in length). Following the Stability-Oracle protocol,
mutations are filtered to enforce at most 30\% sequence identity with the test set, yielding a
reduced set of approximately 117,000 mutation-induced stability changes ($\Delta\Delta G$), referred
to as cDNA117k.\\

(ii) {T2837} (test set): assembled by combining several commonly used benchmarking datasets of
experimentally measured $\Delta\Delta G$ values, including S669~\citep{pancotti_predicting_2022},
myoglobin~\citep{li_predicting_2020}, ssym~\citep{pucci_quantification_2018}, and
p53~\citep{caldararu_base_2021}. The resulting test set comprises 2,837 mutations.


We obtained the cDNA117k and T2837 datasets from the Stability-Oracle~\citep{diaz_stability_2024}
github repository (\url{https://github.com/danny305/StabilityOracle/tree/master}). At the time of
access, the residue indices provided in the dataset did not correspond to the residue numbering in
the original PDB files, but instead to an intermediate, post-processed representation that was not
documented in sufficient detail to allow straightforward recovery of the original numbering. To
ensure consistency with structural data, we manually corrected the residue numbers in the CSV files
to match those in the corresponding PDB structures. The corrected versions of these datasets are
included in our repository.\\



\item {\em ThermoMPNN dataset.} 
ThermoMPNN also leveraged Megascale~\citep{tsuboyama_mega-scale_2023}, but used datasets \# 2 and \#
3, curated and split by the authors at a 25\% sequence-identity cutoff into 216k training and 28k
test mutation effects~\citep{dieckhaus_transfer_2024}. We used the train, valid, and test splits of
the Megascale dataset as curated in ThermoMPNN, and provided on their github repository
(\url{https://github.com/Kuhlman-Lab/ThermoMPNN}). Throughout the text, we refer to these datasets
as the Megascale training (train + valid) and testing sets. We downloaded the structures'
\texttt{.pdb} files from \url{https://zenodo.org/records/7992926}. We note that the Megascale
training set was not controlled for maximum similarity with the T2837 dataset, and six of the T2837
proteins ($\sim 5\%$) have sequence homologs in the Megascale training set with sequence similarity
above 90\%.
\\\\
\end{itemize}

\noindent {\bf Binding effect prediction.} For predicting mutational effects on binding affinity, we
fine-tuned HERMES on the SKEMPI v2.0 dataset~\citep{jankauskaite_skempi_2019}, using 3-fold
cross-validation.

\begin{itemize} 
\item{\em SKEMPI dataset:}
After removing duplicate experimental entries, SKEMPI v2.0 contains 5,713 measurements of binding
free-energy changes ($\Delta\Delta G^{\text{binding}}$) spanning 331 protein–protein complex
structures. We further restrict the dataset to complexes with at least 10 annotated mutations,
resulting in 116 structures and 5,025 mutations. Finally, we retain only single-point mutations,
yielding 93 structures and 3,485 mutations.

SKEMPI provides metadata explicitly designed to support leakage-aware evaluation via two fields:
\emph{hold-out type} and \emph{hold-out proteins} (see SKEMPI documentation:
\url{https://life.bsc.es/pid/skempi2/info/faq_and_help}). Briefly, {\em hold-out type} groups
complexes into broad categories (e.g., protease-inhibitor, antibody-antigen, and pMHC-TCR), while
{\em hold-out proteins} lists PDB identifiers and/or hold-out types that should be co-held-out to
avoid training on closely related binding sites. Using this information, we define three split
strategies:\\


(i) \textit{Easy:} random splitting without using hold-out metadata.\\

(ii) \textit{Medium:} we enforce that all entries linked via a mutation's \emph{hold-out proteins}
annotation are assigned to the same split (but we do not additionally group by \emph{hold-out
type}).\\

(iii) \textit{Difficult:} we enforce that all complexes sharing the same \emph{hold-out type} are
assigned to the same split, producing a stringent evaluation of generalization across interaction
classes.\\

In cases where a complex is associated with multiple hold-out types, we assign it to a single type
by randomly selecting one of the available labels.

\end{itemize}

\subsection{Baseline models}
Here, we describe the baseline models used to benchmark HERMES on mutational effect prediction.\\

\noindent \textbf{ProteinMPNN \citep{dauparas_robust_2022}.} ProteinMPNN is an inverse-folding model
that samples amino-acid sequences conditioned on a protein backbone (optionally with a partial
sequence fixed). Because it also outputs per-site amino-acid probabilities, we used it to score
mutational effects via the log-likelihood ratio in Eq.~\ref{eq:log_ratio_wt_and_mt}. As for HERMES,
we evaluate ProteinMPNN models trained with two noise levels: 0.02 \AA (virtually no noise) and 0.30
\AA. We used, and provide, scripts to infer mutation effects built upon a public fork of the
ProteinMPNN repository (\url{https://github.com/gvisani/ProteinMPNN-copy}).\\
\\
\noindent\textbf{ThermoMPNN~\citep{dieckhaus_transfer_2024}.} ThermoMPNN is a thermodynamic
stability predictor built on top of ProteinMPNN. For a given structure, it extracts the final
ProteinMPNN residue embedding at the site of interest and feeds it to a separate head that predicts
per–amino-acid $\Delta G$ values, from which $\Delta\Delta G$ is computed. Similar to HERMES, this
formulation enforces the permutation symmetry of mutational effects by construction, without
requiring data augmentation. For our experiments, we used native functionalities in the ThermoMPNN
repository (\url{https://github.com/Kuhlman-Lab/ThermoMPNN}). \\
\\
\noindent \textbf{Stability-Oracle~\citep{diaz_stability_2024}.} Similar to HERMES, Stability-Oracle
is trained in two stages. First, a graph-attention network is pre-trained to predict masked amino
acids from their local atomic environment (``neighborhood"). Next, the embeddings from the
pre-trained model is used to regress over mutation-effect. Specifically, For a target site on a
structure, the masked-neighborhood embedding $h$ is extracted from the pre-trained network and
concatenated separately with embeddings of the ``from" and ``to" amino acids. Each concatenated
input is passed through a transformer to produce amino-acid–specific embeddings $e_{aa_{\from}}$ and
$e_{aa_{\tto}}$, whose difference $(e_{aa_{\tto}}-e_{aa_{\from}})$ is fed to a two-layer MLP to
predict $\Delta\Delta G_{aa_{\from}\rightarrow aa_{\tto}}$. This construction is
permutation-symmetric up to the final MLP, since each $e_{aa}$ is computed independently; symmetry
is broken only by the MLP, and would have been preserved by a bias-free linear layer. The original
method enforces symmetry during training via $19\times$ data augmentation. We report performance
scores calculated using predictions provided in the Stability-Oracale's repository
(\url{https://github.com/danny305/StabilityOracle}).\\
\\
\noindent \textbf{RaSP \citep{blaabjerg_rapid_2023}.} Similar to HERMES, RaSP is trained in two
steps. First, a 3D CNN is pre-trained to predict masked amino acids from their local atomic
environment (``neighborhood"). Then, a small fully-connected neural network with a single output is
trained to regress over mutation effects, using as input neighborhoods' embeddings from the 3DCNN,
the one-hot encodings of wildtype and mutant amino-acids, and the wildtype and mutant amino-acids'
frequencies in the pre-training data.
RaSP is fine-tuned on the stability effect of mutations $\Delta\Delta G$, computationally determined
with Rosetta~\citep{chaudhury_pyrosetta_2010}. We report performance scores calculated using
predictions provided in the authors' repository
(\url{https://github.com/KULL-Centre/_2022_ML-ddG-Blaabjerg}).\\
\\
\textbf{RDE-Network~\citep{luo_rotamer_2022}.} RDE stands for Rotamer Density Estimator. This model
consists of first a graph neural network encoder that is trained via a normalizing flow objective to
predict distributions of side-chain conformations. Then, a prediction head is added, and it is
trained on $\Delta\Delta G_{\text{binding}}$ effects from the SKEMPI dataset. We report performance
scores as provided in ref.~\citep{luo_rotamer_2022}.
\\
\\
\textbf{MIF-Network~\citep{luo_rotamer_2022}.} This model's architecture is the same as the encoder
in RDE-Network. It is first pre-trained on the task of wild-type amino acid classification. Then, a
prediction head is added to the encoder, and it is trained on $\Delta\Delta G_{\text{binding}}$
effects from the SKEMPI dataset~\citep{jankauskaite_skempi_2019} in the same manner as RDE-Network.
We report performance scores as provided in Ref.~\citep{luo_rotamer_2022}.\\
\\
\textbf{Pythia-PPI~\citep{tao_reliable_2025}.} Pythia-PPI uses a graph neural network encoder
pre-trained for wild-type amino acid classification, followed by a $\Delta\Delta G$ prediction
module with two heads: one for $\Delta\Delta G_{\text{stability}}$ and one for $\Delta\Delta
G_{\text{binding}}$. The model is fine-tuned jointly on stability labels from
FireProtDB~\citep{stourac_fireprotdb_2021} and binding labels from
SKEMPI~\citep{jankauskaite_skempi_2019}, using a validation selected 20:80 stability:binding mixing
ratio. The resulting checkpoint (``Vanilla Pythia-PPI") is then further trained via
self-distillation by fine-tuning on its own predictions over all SKEMPI complex structures to obtain
the final model (Pythia-PPI). We should note that these two procedures (i.e., joint fine-tuning on
stability and binding with a validation selected mixing ratio and self-distillation) could also be
applied to HERMES to potentially improve performance. We report performance scores as provided in
ref.~\citep{tao_reliable_2025}.


\subsection{ESMFold for computational modeling of protein structures}
For the analysis in Fig.~\ref{fig:radial_plots__esmfold}, we used the ESM Metagenomic Atlas API to
fold each sequence individually (\url{https://esmatlas.com/resources?action=fold}).


\subsection{Structure-conditioned reversibility analysis}

In Figs.~\ref{fig:ssym_figure}~and~\ref{fig:ssym_antisymmetry_reversibility_score_results}, we
characterized structure-conditioned reversibility on the SSym
dataset~\cite{pucci_quantification_2018}, which contains wild-type and single-mutant structures for
352 mutations across 19 proteins~\citep{pancotti_predicting_2022}.

We assess structure-conditioned reversibility by whether the effects of forward $\alpha\to \beta$
and reverse $\beta\to \alpha$ mutations are equal, using the outgoing structural context to compute
mutational effects. Specifically, the forward mutational effect $\Delta \hat
F^{(\text{model})}_{\alpha\to \beta}$ is computed by conditioning on the structure $X_\alpha$, and
the reverse effect $\Delta \hat F^{(\text{model})}_{\beta\to \alpha}$ is computed using $X_\beta$,
\begin{equation}
\begin{aligned}
    \Delta \hat F^{(\text{model})}_{\alpha\to \beta} = \log P(\beta | X_{\alpha}) - \log P(\alpha
    |X_{\alpha}) \\
    \Delta \hat F^{(\text{model})}_{\beta\to \alpha} = \log P(\alpha | X_{\beta}) - \log P(\beta
    |X_{\beta})
\end{aligned}
\end{equation}
With this definition, we do not expect reversibility, as forward and reverse transitions are
conditioned on different structures.\\


\noindent{\bf Reversibility score.} For the analysis in
Figs.~\ref{fig:ssym_figure}~and~\ref{fig:ssym_antisymmetry_reversibility_score_results} we construct
a reversibility score as a mean squared error normalized to be between -1 and 1,
\EQ
\text{rev. score} = 1 - \frac{\text{mean}((\Delta\log p_{fwd}+\Delta\log
p_{rev})^2)}{\text{mean}({\Delta\log p_{fwd}}^2)+\text{mean}({\Delta\log p_{rev}}^2)}
\label{eq:rev_score}
\EE
where $\Delta \log p_{fwd} = \log p(\mt | X_\wt) - \log p(\wt | X_\wt)$ is the forward mutational
effect computed by conditioning on the wild type structure, and $\Delta \log p_{rev} = \log p(\wt |
X_\mt) - \log p(\mt | X_\mt)$ is the reverse mutational effect computed by conditioning on the
mutants structure. The averages (mean) are computed over the 352 mutations in the SSym dataset, or
stratified as needed by the desired criteria.
With this definition, a larger value of {\bf rev. score} implies more degree of reversibility
between forward and reverse mutations. \\


\noindent{\bf Sequence similarity calculation between the Ssym dataset and pre-training proteins.}
We stratified the data based on their similarity to the training data. To do so, we used BLASTp via
NCBI BLAST+~\cite{camacho_blast_2009} with individual chains in the SSym structures as
queries~\cite{pucci_quantification_2018}, and individual chains in the pre-training set as the
database; for each SSym chain, we then considered its maximum similarity to any sequence in the
pre-training set. We found that 9 Ssym proteins have similarity below 70\% to the pre-training
proteins (only 5 are below 50\%, none are below 40\%), and 10 proteins have similarity above 70\%
(comprising the two panels in Fig.~\ref{fig:ssym_figure}B).\\


\subsection{Statistical comparison of model performances on classification metrics}

\noindent{\bf Comparing models' performances on the same tasks.}
Figures~\ref{fig:radial_plots__noise}~and~\ref{fig:megascale_precision_recall_f1_by_size_cutoff_0p0}
report classification performance (precision, recall, F1, etc.) for predicting the stability effects
of mutations across models.
Figures~\ref{fig:t2837_and_megascale_pairwise_pvalues_permutation}~and~\ref{fig:pairwise_pvalues_permutation_bucketed_by_sizereduced_precision_recall_f1}
report pairwise significance tests (p-values) for differences in these metrics using permutation
tests, with the null hypothesis that two models' predictions are exchangeable. To compute these
p-values, for each model pair, we generated permuted prediction sets by swapping paired prediction
between models with probability 0.5, and repeat this procedure for 1000 random seeds. The p-value is
the fraction of permutations in which the metric difference between the permuted sets is at least as
large ($\geq$) as the observed difference. We correct the computed p-values for multiple comparisons
using Holm–Bonferroni across all model pairs within each metric (i.e., within each panel in
Figs.~\ref{fig:t2837_and_megascale_pairwise_pvalues_permutation},~\ref{fig:pairwise_pvalues_permutation_bucketed_by_sizereduced_precision_recall_f1}).\\


\noindent{\bf Comparing a model's performance across tasks.}
Figure~\ref{fig:megascale_precision_recall_f1_by_size_cutoff_0p0} reports classification performance
(precision, recall, F1, etc.) for predicting mutation stability effects across mutation size
categories (tasks).
Figure~\ref{fig:pvalues__bucketed_by_size__between_small_large_buckets__vertical} reports, for each
model, pairwise significance tests for differences in these metrics between tasks. We estimate
p-values via bootstrap resampling: for each task pair, we repeatedly (1000 seeds) draw bootstrap
samples of predictions within each task, recompute the metric difference, and define the p-value as
the fraction of bootstrap replicates in which the resampled metric difference is at least as large
($\geq$) as the observed metric difference. We apply Holm–Bonferroni correction across all task
pairs within each metric.\\



\noindent{\bf Significance of the number of retrieved antigen-stabilizing mutations by different
models.}
In
Figures~\ref{fig:antigen_figure_main_results},~\ref{fig:antigen_figure_of_different_types_of_mutations}
and Table~\ref{table:antigen_results}, we report, for each model, the number of antigen-stabilizing
mutations retrieved ($x$) out of a total of $N=33$ experimentally verified such mutations. Retrieval
is based on amino-acid ranking: a mutation is counted as retrieved if (i) the mutant amino acid is
ranked better than the wild type ($r_\mt<r_\wt$), and (ii) that its rank is below or equal to a
certain threshold $R$, with $R=3$ for ``strongly suggested" and $R=6$ for ``moderately suggested."
We assess significance in the number of retrieved antigen-stabilizing mutations by a given model
with a one-sided binomial test: for each model, the p-value is the probability under a null model of
recovering at least $x$ successes in $N$ trials with per-mutation success probability
$p_{\text{null}}$, i.e., $p = 1-\text{BinomCDF}(x-1, N, p_{\text{null}})$. We adjust p-values for
multiple testing using Holm–Bonferroni, and report the resulting p-values in
Figs.~\ref{fig:antigen_pvalues_vs_random}~and~\ref{fig:antigen_pvalues_vs_blosum62}.\\
To compute the p-values, we consider two null models:
\begin{enumerate}[(i)]
    \item \textbf{Random null.} For each mutation, we sample ranks for the mutant and wild type
    uniformly without replacement: $r_{\mathrm{mt}}\sim \mathrm{Unif}\{1,\dots,20\}$ and
    $r_{\mathrm{wt}}\sim \mathrm{Unif}\{1,\dots,20\}\setminus{r_{\mathrm{mt}}}$. A mutation is a
    ``success" if $r_{\mathrm{mt}}<r_{\mathrm{wt}}$ and $r_{\mathrm{mt}}\leq R$, which occurs with
    probability $p_{null} = \frac{1}{19 \times 20}\sum_{i=1}^{R} (20 - i)$.\\
    \item \textbf{BLOSUM62 null.} We simply set $p_{null} = x_{B62} / N$, where $x_{B62}$ is the
    number of mutations with BLOSUM62 rank $r^{B62}_\mt \leq R$. We note that the rank of the
    wild-type is always 1 for BLOSUM62, so we omit the condition that the rank of the mutant has to
    be lower than the rank of the wild-type, which is instead applied to all other models, including
    the random null model.
\end{enumerate}



\subsection{Using Rosetta to score antigen-stabilizing mutations}
In Figures~\ref{fig:antigen_figure_main_results},~\ref{fig:antigen_figure_of_different_types_of_mutations},
~\ref{fig:antigen_heatmap} and Tables~\ref{table:antigen_results},~\ref{table:speed_on_antigens}, we reported the Rosetta
scores for verified antigen-stabilizing mutations.
\GM{In Figures~\ref{fig:rsv_f_rosetta_hits_main_figure}~and~\ref{fig:denv_e_rosetta_hits_main_figure}
we used Rosetta scores as oracles for mutations suggested by other models.}
To compute these scores, we used the
PyRosetta software~\citep{chaudhury_pyrosetta_2010} to model protein structures, with wild-type
structures serving as template.
For each target sequence containing a single point mutation, we threaded the mutant sequence onto
all chains of the oligomeric template. Each resulting model was then subjected to a high-resolution
refinement protocol using Rosetta's FastRelax application. This protocol involved five cycles of
side-chain rotamer repacking followed by gradient-based energy minimization of backbone ($\phi$,
$\psi$) and side-chain ($\chi$) torsion angles, guided by the \texttt{ref2015\_cart} all-atom energy
function.
\GM{To improve conformational sampling, we repeated the threading-and-relaxation procedure 20
times per mutation for use in the analysis against validated mutations, and 50 when used as an oracle.
For each mutant, we report the mean Rosetta Energy Unit (REU) score of the 20\%
lowest-energy (most favorable) relaxed models (5 for the analysis against validated mutations,
13 when used as an oracle).}



\section*{Supplementary Information}




\subsection*{Extended analysis of antigen stabilization}

Here, we present a more detailed analysis of our antigen stabilization predictions from
Figures~\ref{fig:antigen_figure_main_results},~\ref{fig:antigen_figure_of_different_types_of_mutations}
and Table~\ref{table:antigen_results}. For each antigen, we contextualize the candidate mutations by
their local structural environment and examine the characteristics and physico-chemical features of
the HERMES-predicted top-ranked antigen-stabilizing substitutions (Fig.~\ref{fig:antigen_heatmap}).
This analysis is intended to help practitioners incorporate these models into design workflows and
to motivate quantitative, context-dependent success criteria for real-world stabilization tasks. As
demonstrated in the main text
(Figs.~\ref{fig:antigen_figure_main_results},~\ref{fig:antigen_figure_of_different_types_of_mutations}),
we find that comparing the rank of the wild-type residue to that of putative stabilizing
substitutions is an informative diagnostic of model behavior and performance.\\\\
\noindent\textbf{RSV-F.} 
RSV-F Cav1 stabilitizing mutations S190F and V207L are well-predicted by different HERMES models
(the stability-fine-tuned models trained on +cDNA117k and +Megascale, and
HERMES-\textit{amortized}), whereas Rosetta and ProteinMPNN struggle to identify these mutations. We
therefore focus on these two residues to analyze, post hoc, the structural context underlying
HERMES' predictions, with the goal of clarifying which classes of stabilizing mutations HERMES is
best suited to characterize.


Both S190F and V207L enhance hydrophobic packing in an underpacked region near the trimer apex
(Fig.~\ref{fig:rsvf_si_figure}). V207L is a significantly more anticipated mutation than S190F,
indicated by a positive BLOSUM62 score; this is reflected in HERMES assigning a higher rank to the
wild-type Val207 relative to the mutant Leu (Fig.~\ref{fig:antigen_heatmap} and
Table~\ref{table:antigen_results}). Notably, the highest-ranking residue predicted by HERMES models
at position 207 is Ile (Fig.~\ref{fig:antigen_heatmap}). Inspection of the native wild-type
structure (PDB ID: 4JHW~\citep{mclellan_structure-based_2013}) suggests that an Ile mutation would
pack exceptionally well in the hydrophobic pocket surrounding residue 207, suggesting V207I might
outperform the identified V207L mutation (Fig.~\ref{fig:rsvf_si_figure}A).

At position 190, the top-ranked substitutions from the wild-type Ser are Val or Ile. Structural
examination suggests both residues can be accommodated within the existing 4JHW structure pocket
without requiring the repacking of surrounding residues (Fig.~\ref{fig:rsvf_si_figure}B).
Conversely, the known stabilizing mutation, Phe, appears to slightly overpack the region in the 4JHW
crystal structure (Fig.~\ref{fig:rsvf_si_figure}B). While the HERMES-\textit{fixed} model does not
strongly prioritize Phe, both HERMES stability-fine-tuned models rank Phe third, behind the smaller
Val and Ile.
 This indicates that fine-tuning on $\Delta\Delta G$ datasets enhances implicit reasoning regarding
 local repacking and structural relation possibilities, particularly for larger residues. The
 HERMES-\textit{amortized} model also outperforms the HERMES-\textit{fixed} model, likely due to the
 model's de-emphasis on strict steric constraints (Table~\ref{table:antigen_results}).\\
\\
\textbf{Universal-HA.} 
We observe robust performance in recovering Universal-HA pH switch mutations across all tested
models. All three mutations carry ``unanticipated" BLOSUM62 scores.

    
The mutational-effect heatmaps (Fig.~\ref{fig:antigen_heatmap}) show that, for HA, ProteinMPNN's
preferences are narrowly concentrated, strongly favoring the reported stabilizing substitutions,
whereas HERMES-\textit{amortized}, HERMES-\textit{fixed} (0.50 + Megascale), and Rosetta yield
substantially broader preference profiles. For H355W, all HERMES models correctly indicate that the
wild-type His is disfavored, but they rank other bulky aromatics (Tyr or Phe) above the stabilizing
Trp.
 Despite being a known stabilizing mutation, Trp does not appear to fit in the 7VDF crystal
 structure without clashes (Fig.~\ref{fig:ha_si_figure}). This discrepancy may stem from the
 limitations of rigid-body mutagenesis (e.g., PyMOL's wizard), as subtle backbone movements are
 likely required to accommodate the mutation. This finding cautions against over-interpreting
 apparent clashes in a single experimentally determined static crystal structure, especially for
 bulky substitutions, and motivates follow-up evaluation with relaxation-aware structure prediction
 models or physics-based repacking/scoring tools to capture local flexibility and repacking that may
 render seemingly sterically forbidden mutations feasible.\\\\
\noindent \textbf{hMPV-F.} 
We first analyze stabilizing mutations in the M-104 hMPV variant~\citep{gonzalez_general_2024}.
The A159L substitution requires the nearby Ile-137 to adopt an alternative
rotamer~\citep{gonzalez_general_2024}. The two $\Delta\Delta G$-fine-tuned HERMES models prioritize
Val (V) and Ile (I) over Leu (L) at position 159. We note that mutation to either Val or Ile at
position 159 would require both Ile-137 and Leu-141 to adopt different rotamers to accommodate their
beta-branched side chains (Fig.~\ref{fig:hmpvf_si_figure}C). Nonetheless, A159V and A159I remain
plausible stabilizing mutations and therefore, warrant experimental screening.
ProteinMPNN identifies Ala as the most favorable residue at this position, indicating limited
sensitivity to larger hydrophobic substitutions. By comparison, the $\Delta\Delta G$-fine-tuned
HERMES models rank the wild-type Ala substantially lower (ordinal rank 5--7), while consistently
favoring larger hydrophobic residues. Low rankings for buried hydrophobic positions may therefore
indicate suboptimal core packing, particularly when alternative residues with greater side-chain
volume are predicted. Although Leu is not the top-ranked substitution, the $\Delta\Delta
G$-fine-tuned HERMES models correctly capture the preference for a residue larger than Ala to
improve packing within this pocket.

The V203I substitution was recovered by nearly all models. A Val to Ile substitution is highly
anticipated, with a BLOSUM62 score of 3. Notably, the $\Delta\Delta G$-fine-tuned HERMES models rank
Ile as more favorable than the wild-type Val, whereas ProteinMPNN, ThermoMPNN,
HERMES-\textit{fixed}, and HERMES-\textit{amortized} prefer the wild-type Val. In accordance with
this mutation being more highly anticipated, it is likely that the pocket occupied by V203 is less
underpacked relative to A159L in the M-104 variant, with the wild-type residue still being quite
highly-favored. The subtle nature of Val to Ile mutations may result in less signal overall when
comparing amino acid ranks as we do in this work.


The V449D substitution replaces a surface-exposed hydrophobic residue with a polar side chain and
introduces hydrogen-bonding contacts to Asn298 (Fig.~\ref{fig:hmpvf_si_figure}A). ProteinMPNN,
ThermoMPNN, and all HERMES variants instead rank Glu as the top substitution; structural inspection
suggests that Glu could form similar hydrogen bonds, and may therefore also be stabilizing
(Fig.~\ref{fig:hmpvf_si_figure}A). Overall, the models performed well in identifying mutations that
enable additional hydrogen bonding.


For V430Q, the predicted $\Delta\Delta G$ reported in the original study is small in magnitude
(albeit negative), suggesting that it should not be classified as strongly stabilizing, as suggested
in the original study~\citep{gonzalez_general_2024}. Consistent with this, neither of the
$\Delta\Delta G$–fine-tuned HERMES models prioritize V430Q. Interestingly, the only model that ranks
V430Q as the top substitution is HERMES-\textit{fixed}.


In the MPV-2c variant~\citep{bakkers_efficacious_2024}, recovery of V112R depends on providing the
model with the trimeric PDB. The introduced Arg forms interprotomer van der Waals contacts and
intraprotomer hydrogen bonds that stabilize the prefusion trimer (Fig.~\ref{fig:hmpvf_si_figure}B).
Because this mechanism is inherently multimeric and not captured by a strictly local energetic
proxy, it may explain why the $\Delta\Delta G$-fine-tuned HERMES models do not prioritize V112R. In
contrast, HERMES-\textit{amortized} ranks Arg second, slightly favoring the more conservative Leu,
which could plausibly adjust hydrophobic interprotomer contacts.


D209E is correctly predicted by ProteinMPNN, HERMES-\textit{fixed}, and HERMES-\textit{amortized},
whereas the $\Delta\Delta G$-fine-tuned HERMES models perform poorly in proposing this mutation. In
particular, the $\Delta\Delta G$-fine-tuned models rank bulky hydrophobics (Leu, Met, Ile, Val)
ahead of the charge-conserving Glu. Although the charged substitution D209E plausibly stabilizes the
prefusion state via favorable polar interactions, the stabilizing potential of these alternative
hydrophobic substitutions remains untested.



E453P replaces a glutamate with proline, a substitution that can be strongly stabilizing by
rigidifying locally flexible regions. Among the evaluated methods, only HERMES-\textit{amortized}
predicts this mutation. ProteinMPNN ranks the wild-type Glu as most favored and Pro as most
disfavored (rank 1 vs. 20), while ThermoMPNN partially recovers the substitution but still favors
the native Glu. Consistent with these mixed signals, the original
study~\citep{bakkers_efficacious_2024} reports that E453P improves stability but reduces expression,
and that E453Q yields weaker stabilization with improved expression. This highlights the poorly
understood coupling between stability and expression and suggests that some models may
systematically downweigh proline substitutions, potentially because prolines are underrepresented in
natural sequences, and they appropriately learned to rarely recommend them.
The fact that HERMES-\textit{amortized}—but not HERMES-\textit{fixed}—recovers E453P suggests that
modeling relaxed neighborhoods helps identify sites that can accommodate Pro's constrained
Ramachandran geometry. Finally, we suspect that the utility and stabilizing ability of proline
substitutions in antigen design is not fully captured by HERMES when fine-tuned with the cDNA117k or
Megascale datasets, as the proteins represented in those datasets are much smaller, and lack the
interplay of conformational switching between prefusion and post-fusion.\\\\
\textbf{DENV-E.} 
We observe mixed performance on DENV-E stabilizing mutations. S29K is recovered only by ProteinMPNN
and ThermoMPNN. Although S29K, T33V, and A35M (``PM4" mutations as a group) are spatially proximal
and may act synergistically, they were originally identified by Rosetta site-saturation mutagenesis,
suggesting they should be accessible to single point-mutation scoring
schemes~\citep{phan_conserved_2022}. In the 1OAN structure~\citep{modis_ligand-binding_2003},
introducing Lys at position 29 appears to create substantial steric clashes across rotamers
(Fig.~\ref{fig:denve_si_figure}A). S29K and A35M are not anticipated by BLOSUM62, whereas T33V is
neutral according to BLOSUM62 (score of $0$, Table~\ref{table:antigen_results}) but is strongly
preferred by all models tested (Fig.~\ref{fig:antigen_figure_main_results} and
Table~\ref{table:antigen_results}). Structurally, T33 lines a hydrophobic pocket, making the
isosteric Val substitution easier to predict. Notably, the PM4 mutations are not present in the
best-available SC12 structure (PDB: 6WY1)~\citep{kudlacek_designed_2021}.


A major challenge in stabilizing DENV-E is strengthening homodimer interactions. A259W and T262R
were mutations made at the dimer interface with the intention of enhancing the strength of the
homodimer. A259W is highly unanticipated (BLOSUM62 = $-3$) but, together with T262R, can form a
favorable cation--$\pi$ interaction while improving packing against the neighboring protomer
(Fig.~\ref{fig:antigen_figure_of_different_types_of_mutations}A.4). Notably, the HERMES model
fine-tuned on Megascale favors A259W despite its apparent synergy with T262R; this could reflect
sensitivity to the interprotomer packing geometry, although a chance effect cannot be excluded. We
also expect the generally weak recovery of T262R across models to improve when Trp is present at
position 259, consistent with the coupled nature of this interface motif.



Like E453P in hMPV-F, the proline substitution T280P is recovered only by HERMES-\textit{amortized},
further highlighting the model's ability to reason about proline accommodation. F279W fills an
underpacked region. However, comparison of the stabilized 6WY1 structure with the native 1OAN
crystal structure suggests that backbone movement (residues 269--281) is required to optimally fit
Trp, potentially aided by T280P (Fig.~\ref{fig:denve_si_figure}B). ThermoMPNN correctly reasons
regarding the Trp introduction, and HERMES-\textit{fixed} $+$ Megascale ranks Trp second. Other
stability-fine-tuned HERMES models prefer Phe, Ile, or Leu, likely adhering more strictly to the
steric limitations of the 1OAN backbone.

We observe strong reasoning across almost all HERMES models, ProteinMPNN, and Rosetta in identifying
G106D, which introduces stabilizing polar and electrostatic contacts
(Fig.~\ref{fig:denve_si_figure}C). Curiously, ThermoMPNN struggles with this specific mutation.\\
\\
\textbf{SARS-CoV-2 Spike.} 
The spike protein of SARS-CoV-2 was initially stabilized through the introduction of two proline
mutations (``S-2P"), and subsequent work identified four additional prolines that further improve
stability and expression of the full-length spike~\citep{hsieh_structure-based_2020}. Here we
evaluate recovery of these four additional proline sites (excluding the original S-2P mutations).
Consistent with its strong performance on proline substitutions, HERMES-\textit{amortized} ranks Pro
as the top choice at all four positions. ProteinMPNN also performs well, recovering Pro at 3/4
sites. \\\\
\textbf{Summary.} 
Our analysis across multiple antigens reveals several consistent themes regarding model performance
and decision-making logic.

First, we observe distinct behaviors regarding hydrophobic packing. The HERMES models fine-tuned on
stability datasets consistently demonstrate enhanced reasoning for hydrophobic core packing. These
models frequently suggest larger hydrophobic residues (e.g., Ile or Phe) to fill underpacked
cavities where models without explicit stability training might prefer the native residue or smaller
conservative substitutions. However, this sensitivity can sometimes lead to ``confusion" among
similar hydrophobic amino acids (e.g., Val vs. Ile vs. Leu), where the models correctly identify the
chemical property needed but may rank several hydrophobic options similarly.

Second, the interplay between structural rigidity and model flexibility is evident in the prediction
of proline mutations. HERMES-\textit{amortized} consistently outperforms other models in identifying
stabilizing proline substitutions. This suggests that the model's training, which incorporates
relaxed neighborhoods, allows it to better identify backbone locations capable of accommodating the
steric constraints of proline, whereas other methods more often penalize these stabilizing mutations
due to perceived steric incompatibilities of the provided protein structure.

Third, we note differences in the mutational landscape profiles predicted by different architectures
(Fig.~\ref{fig:antigen_heatmap}. ProteinMPNN tends to produce narrow substitution profiles, often
assigning very low probabilities to non-native residues unless the signal is overwhelmingly strong.
In contrast, HERMES models--particularly HERMES-\textit{amortized} and those fine-tuned on Megascale
stability data--exhibit broader mutational profiles. This broader landscape may be more advantageous
for design applications, as it provides a richer set of plausible candidate hypotheses for
experimental validation and may guide a practitioner's intuition regarding the nature or predicted
effects of various substitutions in a more granular manner.

Finally, while ProteinMPNN and Rosetta remain powerful tools for sequence recovery, the
HERMES-\textit{amortized} and stability-fine-tuned models demonstrate a unique capacity to
prioritize mutations that improve local packing density even when such mutations appear sterically
challenging in the provided structure. This highlights the utility of using an ensemble of models to
capture different modes of stabilization, from electrostatic optimization to hydrophobic core
repacking and backbone rigidification.


